\documentclass[11pt]{article}

%\documentclass[17pt]{extarticle}

%\usepackage{extsizes}
\usepackage{indentfirst}
\usepackage[utf8x]{inputenc}
\usepackage[T1]{fontenc}
\usepackage[english,lithuanian]{babel}
\usepackage{array}
\usepackage{caption}

\usepackage{amsmath, amsthm, amssymb}
\usepackage{graphicx}
\usepackage{setspace}
\usepackage{verbatim}
\usepackage[left=3cm,top=2cm,right=1.5cm,bottom=2cm]{geometry}
 %\textwidth 6.5in
 %\textheight 9.00in

  \onehalfspacing

% Author's MACROS
 \newcommand{\EE}{\mathbb{E}\,} % Mean
 \newcommand{\ee}{{\mathrm e}}  % nice exponent
 \newcommand{\dd}{{\mathrm d}}
 \newcommand{\RR}{\mathbb{R}}
 %Macros end

\begin{document}
\selectlanguage{lithuanian}

\begin{titlepage}
\vskip 20pt
\begin{center}
\includegraphics[scale=0.5]{MIF}
\end{center}

%%%%%%%%%%%%%%%%%%%%%%%%%%%%%%%%%%%%%%%%
% TITULINIO PUSLAPIO TEKSTAS
%%%%%%%%%%%%%%%%%%%%%%%%%%%%%%%%%%%%%%%%

\vskip 20pt
\centerline{\bf \large \textbf{VILNIAUS UNIVERSITETAS}}
\bigskip
\centerline{\large \textbf{MATEMATIKOS IR INFORMATIKOS FAKULTETAS}}
\bigskip
\centerline{\large \textbf{BIOINFORMATIKOS BAKALAURO STUDIJŲ PROGRAMA}}

\vskip 90pt
\begin{center}
    {\bf \LARGE \emph{tbx5} ir \emph{tcf21} genų įtakos regeneracijai tyrimai}
\end{center}
\begin{center}
    {\bf \Large Research of the \emph{tbx5} and \emph{tcf21} genes' functions in regeneration}
\end{center}
\vskip 20pt
\centerline{\bf \Large \textbf{Kursinis darbas}}
\bigskip
\vskip 50pt

\hskip 140pt {\Large Autorius: Danielė Stasiūnaitė}

\hskip 140pt{\Large VU el. p.: (daniele.stasiunaite@mif.stud.vu.lt)}
\bigskip
\vskip 20pt

\hskip 140pt {\Large Darbo vadovas: J. Asist. Kotryna Kvederavičiūtė}
\vskip 60pt
\vskip 60pt
\centerline{\large \textbf{Vilnius}}
\centerline{\large \textbf{2022}}
\newpage
\end{titlepage}

\selectlanguage{lithuanian}

%%%%%%%%%%%%%%%%%%%%%%%%%%%%%%%%%%%%%%%%
% TURINIO PUSLAPIS
%%%%%%%%%%%%%%%%%%%%%%%%%%%%%%%%%%%%%%%%  
\tableofcontents
\newpage

%%%%%%%%%%%%%%%%%%%%%%%%%%%%%%%%%%%%%%%%
% LIETUVIŠKOS SANTRAUKOS PUSLAPIS
%%%%%%%%%%%%%%%%%%%%%%%%%%%%%%%%%%%%%%%%  
\section*{Santrauka}
Darbo santrauka.\\

\textbf{Raktiniai žodžiai: TF, ChIP-seq}
\newpage

%%%%%%%%%%%%%%%%%%%%%%%%%%%%%%%%%%%%%%%%
% ANGLIŠKOS SANTRAUKOS PUSLAPIS
%%%%%%%%%%%%%%%%%%%%%%%%%%%%%%%%%%%%%%%%
\section*{Summary}
Short summary of results.\\

\textbf{Keywords?}
\newpage

%%%%%%%%%%%%%%%%%%%%%%%%%%%%%%%%%%%%%%%%
% ĮVADO PUSLAPIS
%%%%%%%%%%%%%%%%%%%%%%%%%%%%%%%%%%%%%%%%
\section{Įvadas}
\newpage

%%%%%%%%%%%%%%%%%%%%%%%%%%%%%%%%%%%%%%%%
% DUOMENŲ APŽVALGA
%%%%%%%%%%%%%%%%%%%%%%%%%%%%%%%%%%%%%%%%
\section{Duomenys}

Tyrime naudoti duomenys atsisiųsti iš \emph{GTRD} (Gene Transcription
Regulation Database) duomenų bazės, saugančios informaciją apie įvairių
transkripcijos faktorių jungimosi prie DNR sekų regionus. Ši duomenų
bazė pasirinkta dėl sistemiškai surinktų ChIP-seq eksperimentų, kurių
metu gauti rezultatai unifikuotai apdorojami ir paruošiami tyrėjų
analizėms. Taip pat \emph{GTRD} duomenų bazėje duomenys saugomi
binariniu anotacijų formatu \emph{bigBed}, leidžiandžiu atvaizduoti
pasirinktą chromosomos regioną genomų naršyklėje sparčiau ir 
efektyviau nei tekstinis \emph{BED} formatas.

Žemiau esančioje lentelėje pateikta informacija apie tyrimui atlikti
naudotus duomenis, surinktus iš naminės pelės \emph{(lot. Mus musculus)}
ląstelių. Visi duomenys gauti iš eksperimentų, kurių metu buvo
siekiama išsiaiškinti \emph{tbx5} ir kitų transkripcijos faktorių
įtaką naujų širdies ląstelių, kardiomiocitų, susidarymui bei širdies
funkcijų atsistatymui po ištikusio infarkto arba kitų, su širdies
pažeidimais susijusių, sutrikimų.

\begin{table}[htb]
    \newcolumntype{M}[1]{>{\centering\arraybackslash}m{#1}}
    \small
    \caption*{1 lentelė. Mėginių charakteristikos}
    \begin{tabular}{|M{2cm}|M{3.5cm}|M{3.7cm}|M{3cm}|M{3cm}|M{0.5cm}|}
    \hline
    \textbf{GTRD ID} & \textbf{Ląstelių tipas} &
                    \textbf{Sąlygos pritaikymas} & \textbf{Antikūnai} &
                    \textbf{PubMed ID}\\
    \hline
    EXP030893 & HL - 1 (širdies raumens) & Kontrolė & - & 21415370\\
    \hline
    EXP058852 & Širdies prieširdžių & Kamienas: C57BL/6; vystymosi
                stadija: p4 & Tbx5 (sc-17866) & 31080136\\
    \hline
    EXP062056 & Pelių naujagimių širdies fibroblastų, ekspresuojančių didelį
                kiekį T antigeno, linija & Kamienas: CD1; sąlyga: sb431542,
                    xav939 & anti-TBX5 (sc-17866x) & 31271750\\
    \hline
    EXP058843 & MEF (embrionų fibroblastai) & Kamienas: C57BL/6; sąlyga: AGHMT
                (2 d.) & anti-Tbx5 (sc-17866) & 31080136\\
    \hline
    EXP058847 & MEF (embrionų fibroblastai) & Kamienas: C57BL/6; sąlyga: GHMT
                (2 d.) & Tbx5 (sc-17866) & 31080136\\
    \hline
    EXP058850 & MEF (embrionų fibroblastai) & Kamienas: C57BL/6; sąlyga: GMT
                (2 d.) & Tbx5 (sc-17866) & 31080136\\
    \hline
    EXP058856 & MEF (embrionų fibroblastai) & Kamienas: C57BL/6; sąlyga: vienas
                faktorius (2 d.) & Tbx5 (sc-17866) & 31080136\\
    \hline
    \end{tabular}
\end{table}

Tyrimui atlikti pasirinkti septynių eksperimentų duomenys tarp kurių buvo
tiriami keturi embrionų fibroblastų kamienai, kuriems dvi dienas taikytos
skirtingos sąlygos: veikimas AGHMT (AKT1 - serino/treonino kinazė 1; GATA4,
HAND2, MEF2C, TBX5 - kardiogeniniai transkripcijos faktoriai), GHMT, GMT
ir tik vienu transkripcijos faktoriumi, kuris nebuvo specifikuotas.

Taip pat tirti pelių naujagimių fibroblastai, veikti inhibitoriais: 
sb431542, skatinančiu kardiomiocitų diferenciaciją iš pliuripotentinių
kamieninių ląstelių, ir xav939, sukeliančiu progenitorinių ląstelių
kardiomiogenezę.

Į duomenų rinkinį įtraukti du mėginiai, kuriems nebuvo taikytos papildomos
sąlygos: kontrolinis HL - 1 širdies raumens ląstelių mėginys bei p4
vystymosi stadijos širdies prieširdžių ląstelių mėginys.
\newpage

%%%%%%%%%%%%%%%%%%%%%%%%%%%%%%%%%%%%%%%%
% METODAI
%%%%%%%%%%%%%%%%%%%%%%%%%%%%%%%%%%%%%%%%

\section{Tyrimo metodai}
\emph{Tbx5} transkripcijos faktoriaus regionų tyrimo analizė atlikta
su R programavimo kalba (4.2.0 versija). Žemiau esančiuose atskiruose
analizės etapų skyriuose nurodomos R bibliotekos ir papildomi komandinės
eilutės įrankiai, kuriais naudojantis pasiekti tarpiniai analizės etapų
rezultatai.

\subsection{\emph{Tbx5} regionų skaičiaus nustatymas mėginiuose}
\emph{Tbx5} regionų skaičius skirtinguose mėginiuose apskaičiuotas
\emph{rtracklayer} bibliotekos sukurtam GRanges objektui, aprašančiam
genomines pozicijas bei su jomis susijusias anotacijas, pritaikius
standartinę R ilgio funkciją \emph{length()}.
Regionų skaičių mėginiuose atvaizduojanti stulpelinė diagrama sukurta
su \emph{ggplot2} bibliotekos \emph{geom\_bar()} funkcija.

\subsection{\emph{Tbx5} regionų skaičiaus pasiskirstymas chromosomose}
Transkripcijos faktoriaus regionų skaičius skirtingose chromosomose
suskaičiuotas, naudojantis standartine R funkcija \emph{length()},
pritaikyta atskiroms chromosomoms, kurių pozicijos aprašytos GRanges
objekte.
Kiekvieno mėginio \emph{Tbx5} transkripcijos faktoriaus pasiskirstymas
chromosomose atvaizduotas su \emph{ggplot2} biblioteka, panaudojus
papildomą funkciją \emph{facet\_wrap()}, sukuriančia grafikus visoms
chromosomoms.

\subsection{Persidengiančių regionų mėginiuose analizė}
Tekstas.

\subsection{\emph{Tbx5} motyvo nustatymas}
Tekstas.

\subsection{Motyvų paieška \emph{de novo}}
Tekstas.



\newpage

\section{Antros dalies pavadinimas}
\subsection{Skyriaus pavadinimas}
Grafikams įterpti:
\begin{figure}[h]

\end{figure}


\subsection{Skyriaus pavadinimas}

\newpage

%%%%%%%%%%%%%%%%%%%%%%%%%%%%%%%%%%%%%%%%
% GAUTŲ REZULTATŲ APŽVALGA
%%%%%%%%%%%%%%%%%%%%%%%%%%%%%%%%%%%%%%%%

\section{Rezultatų apžvalga}
Pateikiami svarbiausi rezultatai, apibendrinimai.

\newpage

%%%%%%%%%%%%%%%%%%%%%%%%%%%%%%%%%%%%%%%%
% LITERATŪROS ŠALTINIAI
%%%%%%%%%%%%%%%%%%%%%%%%%%%%%%%%%%%%%%%%

\bibliographystyle{plain}
\begin{thebibliography}{99}

\bibitem{YLS16}1. Yang Yang; R. Leipus, J. \v{S}iaulys. Asymptotics for randomly weighted and stopped dependent sums, \emph{Stochastics: an international journal of probability and stochastic processes}, 2016, \textbf{88}(2), p.p. 300-319.
 \bibitem{E05} T. Erhardsson. Stein’s method for Poisson and compound Poisson approximation, \emph{In: An Introduction to Stein’s Method. Lect. Notes Ser. Inst. Math. Sci. Natl. Univ. Singap. v. 4}, Singapore: Singapore Univ. Press, 2005, p.p. 61–113.
\bibitem{KGDD08} R. Kaas, M. Goovaerts, J. Dhaene  and M. Denuit. \emph{Modern Actuarial Risk Theory:
using R.} (Secon. edt.), Springer-Verlag, Berlin, Heidelberg,  2008, 393 p.
\bibitem{St12} V. Stakėnas, \emph{Tikimybių teorija ir matematinė statistika} (rankraštis), 2012, 178 p. , http://www.statistika.mif.vu.lt/atsisiuntimui/statistika/

\end{thebibliography}
\newpage

%%%%%%%%%%%%%%%%%%%%%%%%%%%%%%%%%%%%%%%%
% PRIEDAI
%%%%%%%%%%%%%%%%%%%%%%%%%%%%%%%%%%%%%%%%

\section{Priedas}
Programiniam kodui pateikti galima naudoti  \verb|\verbatim| komandą:
\small  %mažesnės raidės
\begin{verbatim}

R arba Python programinis kodas.

\end{verbatim}



\section{Literatūros apžvalga}
Literatūros darbo tema apžvalga. Cituojant \cite{YLS16,KGDD08}. Nepamirškite pdflatex įvykdyti bent du kartus, kad pamatytumėte citavimą, ne klaustukus.


\end{document}
